% 分类目录：dl
\documentclass[12pt, a4paper]{article}
\usepackage[margin=2.5cm]{geometry}
\usepackage{ctex}
\usepackage{amsmath, amssymb}
\usepackage{listings}
\usepackage{xcolor}
\usepackage{tabularx}
\usepackage{hyperref}

\lstset{
    language=Python,
    basicstyle=\ttfamily\small,
    keywordstyle=\color{blue},
    commentstyle=\color{green!50!black},
    stringstyle=\color{red},
    showstringspaces=false,
    breaklines=true,
    numbers=left,
    numberstyle=\tiny\color{gray},
    frame=single
}

\title{知识点：超参数优化 (Hyperparameter Optimization)}
\author{}
\date{}

\begin{document}
\maketitle

\section{核心定义}
\textbf{中文定义}：超参数优化（HPO）是寻找一组最优的算法配置（如学习率、网络层数、Batch Size等）的过程。与模型参数（通过梯度下降学习）不同，超参数是在训练阶段开始前手动设置的，且无法直接通过目标函数求导来更新。

\textbf{英文定义}：Hyperparameter Optimization (HPO) is the process of finding the optimal configuration for an algorithm (e.g., learning rate, number of layers, batch size). Unlike parameters, which are learned via gradient descent, hyperparameters are set before training and cannot be directly updated through gradient-based optimization of the objective function.

\section{传统搜索方法}

\subsection{网格搜索 (Grid Search)}
\textbf{方法}：穷举搜索预定义的参数组合。
\\ \textbf{局限}：搜索空间随参数数量指数级增长（维度灾难），且容易在非核心方向浪费大量资源。

\subsection{随机搜索 (Random Search)}
\textbf{原理}：在指定范围内随机采样。
\\ \textbf{优势}：研究证明，随机搜索通常比网格搜索更高效，因为它能更灵活地在不同维度的连续空间中探索，特别是当某些参数对结果影响远大于其他参数时。

\section{启发式与黑盒优化}

\subsection{贝叶斯优化 (Bayesian Optimization)}
\textbf{核心思想}：构建一个“代理模型”（Surrogate Model）来模拟真实的损失函数，并通过“采集函数”（Acquisition Function）决定下一个搜索点。
\begin{itemize}
    \item \textbf{代理模型}：常用高斯过程 (GP) 或 TPE (Tree-structured Parzen Estimator)。
    \item \textbf{采集函数}：权衡“利用”（在已知好的区域精搜）与“探索”（在未知的区域广搜），如 EI (Expected Improvement)。
\end{itemize}

\section{基于多臂老虎机的优化 (Bandit-based)}

\subsection{Successive Halving (连续减半)}
在训练初期同时运行多个配置，定期淘汰表现最差的一半，将资源集中在有潜力的配置上。

\subsection{Hyperband}
\textbf{改进}：解决了 Successive Halving 对资源分配初始点的依赖。通过动态平衡实验数量和每个实验的训练步数，在同样的计算资源下达到更好的效果。

\subsection{BOHB}
将 \textbf{贝叶斯优化} 与 \textbf{Hyperband} 结合。利用贝叶斯模型指导 Hyperband 的采样过程，是目前 AutoML 领域非常高效的方案。

\section{常用 HPO 框架对比}
\begin{table}[h]
\centering
\begin{tabularx}{\textwidth}{|l|X|X|}
\hline
\textbf{框架} & \textbf{核心特点} & \textbf{适用场景} \\
\hline
Optuna & 定义即运行（Define-by-Run）、高效采样 & 研究、日常调参、支持各种库 \\
\hline
Ray Tune & 分布式支持、可扩展性极强 & 大规模工业级训练、多节点搜索 \\
\hline
Scikit-Optimize & 简单易用、基于 GP & 传统 ML、小规模搜索 \\
\hline
\end{tabularx}
\end{table}

\section{关键代码片段 (Optuna)}
\begin{lstlisting}
import optuna

def objective(trial):
    # 1. 定义搜索空间
    lr = trial.suggest_float("lr", 1e-5, 1e-1, log=True)
    num_layers = trial.suggest_int("num_layers", 1, 3)
    
    # 2. 模拟训练与评估
    accuracy = train_and_eval(lr, num_layers)
    return accuracy

# 3. 创建研究对象并运行
study = optuna.create_study(direction="maximize")
study.optimize(objective, n_trials=100)

print(f"Best params: {study.best_params}")
\end{lstlisting}

\section{深度面试题}

\textbf{问：对于学习率这类参数，为什么通常建议在对数空间进行搜索？}\\
答：因为学习率对模型的影响是乘性（Scale-wise）的而非加性的。在 $[0.001, 0.01]$ 之间的变化对模型的影响程度通常与 $[0.01, 0.1]$ 相当。线性搜索会导致在小数值区域探索严重不足。

\textbf{问：什么是早停 (Early Stopping) 在 HPO 中的作用？}\\
答：它是一种资源节省机制。如果某个参数组合在训练前几个 Epoch 的表现明显差于当前最优，HPO 引擎可以提前剪枝（Pruning）该实验，从而将节省下来的计算资源分配给更有潜力的组合。

\section{参考文献}
\begin{enumerate}
    \item Bergstra, J., \& Bengio, Y. (2012). \textit{Random Search for Hyper-Parameter Optimization}. Journal of Machine Learning Research.
    \item Li, L., et al. (2017). \textit{Hyperband: A Novel Bandit-Based Approach to Hyperparameter Optimization}. JMLR.
    \item Akiba, T., et al. (2019). \textit{Optuna: A Next-generation Hyperparameter Optimization Framework}. KDD.
\end{enumerate}

\end{document}
